% Options for packages loaded elsewhere
\PassOptionsToPackage{unicode}{hyperref}
\PassOptionsToPackage{hyphens}{url}
\documentclass[
]{article}
\usepackage{xcolor}
\usepackage{amsmath,amssymb}
\setcounter{secnumdepth}{-\maxdimen} % remove section numbering
\usepackage{iftex}
\ifPDFTeX
  \usepackage[T1]{fontenc}
  \usepackage[utf8]{inputenc}
  \usepackage{textcomp} % provide euro and other symbols
\else % if luatex or xetex
  \usepackage{unicode-math} % this also loads fontspec
  \defaultfontfeatures{Scale=MatchLowercase}
  \defaultfontfeatures[\rmfamily]{Ligatures=TeX,Scale=1}
\fi
\usepackage{lmodern}
\ifPDFTeX\else
  % xetex/luatex font selection
\fi
% Use upquote if available, for straight quotes in verbatim environments
\IfFileExists{upquote.sty}{\usepackage{upquote}}{}
\IfFileExists{microtype.sty}{% use microtype if available
  \usepackage[]{microtype}
  \UseMicrotypeSet[protrusion]{basicmath} % disable protrusion for tt fonts
}{}
\makeatletter
\@ifundefined{KOMAClassName}{% if non-KOMA class
  \IfFileExists{parskip.sty}{%
    \usepackage{parskip}
  }{% else
    \setlength{\parindent}{0pt}
    \setlength{\parskip}{6pt plus 2pt minus 1pt}}
}{% if KOMA class
  \KOMAoptions{parskip=half}}
\makeatother
\usepackage{longtable,booktabs,array}
\newcounter{none} % for unnumbered tables
\usepackage{calc} % for calculating minipage widths
% Correct order of tables after \paragraph or \subparagraph
\usepackage{etoolbox}
\makeatletter
\patchcmd\longtable{\par}{\if@noskipsec\mbox{}\fi\par}{}{}
\makeatother
% Allow footnotes in longtable head/foot
\IfFileExists{footnotehyper.sty}{\usepackage{footnotehyper}}{\usepackage{footnote}}
\makesavenoteenv{longtable}
\setlength{\emergencystretch}{3em} % prevent overfull lines
\providecommand{\tightlist}{%
  \setlength{\itemsep}{0pt}\setlength{\parskip}{0pt}}
\usepackage{bookmark}
\IfFileExists{xurl.sty}{\usepackage{xurl}}{} % add URL line breaks if available
\urlstyle{same}
\hypersetup{
  hidelinks,
  pdfcreator={LaTeX via pandoc}}

\author{}
\date{}

\begin{document}

\section{\texorpdfstring{The Horizon-Condensation Law
(\(\mathcal{C}_\tau = 1/2\pi\)): Speed, Temperature, and Weyl-Invariant
Thermodynamics Across Physical
Regimes}{The Horizon-Condensation Law (\textbackslash mathcal\{C\}\_\textbackslash tau = 1/2\textbackslash pi): Speed, Temperature, and Weyl-Invariant Thermodynamics Across Physical Regimes}}\label{the-horizon-condensation-law-mathcalc_tau-12pi-speed-temperature-and-weyl-invariant-thermodynamics-across-physical-regimes}

\textbf{Author:} Luís Pimenta Damásio\\
\textbf{ORCID:}
\href{https://orcid.org/0009-0001-7606-7557}{0009-0001-7606-7557}\\
\textbf{Affiliation:} Independent Researcher, Lisbon, Portugal\\
\textbf{Date:} November 22, 2026\\
\textbf{Series:} Theory of Objective Time (\(\text{AEON-}T_{\mu\nu}\))
Research Series --- Paper 9\\
\textbf{License:} Creative Commons Attribution 4.0 International (CC BY
4.0)

\begin{center}\rule{0.5\linewidth}{0.5pt}\end{center}

\subsection{Abstract}\label{abstract}

In standard theoretical physics, kinematics
(\(v = \Delta x / \Delta t\)) and thermodynamics (\(T\)) are formulated
as derivative statistical or coordinate constructions laid atop an
assumed four-dimensional spacetime background (\(\mathbb{R}^4\),
\(g_{\mu\nu}\)). Within the Theory of Objective Time
(\(\text{AEON-}T_{\mu\nu}\)), where Objective Time (\(\rho_\tau > 0\))
and ordinal persistence (\(A_s^v \to A_{s'}^{v+1}\)) constitute the
primary ontological substrate, both \textbf{Speed} (\(v/c\)) and
\textbf{Temperature} (\(T\)) emerge directly as the two orthogonal modes
through which bounded, open-system physical clocks (\emph{Centers of
Resolution}) expend their finite internal cadence budget. Speed
represents \emph{coherent, directional spatial re-situating} paid out of
internal clock ticks (\(dt_{\mathrm{pre}} = dt/\gamma\)), whereas
temperature represents the \emph{incoherent, multi-node rate of
Two-Stage Condensations} (\(\Phi_\tau\)) driven by boundary coupling
(\(\Delta\Omega > 0\)).

By unifying these modes in the \textbf{Cadence Partition Identity}
across the \textbf{Weyl Scale Elevator} and analyzing horizon truncation
under sustained acceleration
(\(\mathcal{L}_{\mathrm{sim}}(a) = c^2/a\)), we derive an exact,
dimensionless, Weyl- and Lorentz-invariant constant governing
open-system phase-locking: \textbf{The Horizon-Condensation Law},
\(\mathcal{C}_\tau \equiv \frac{k_B T \cdot \mathcal{L}_{\mathrm{sim}}}{\hbar c} = \frac{1}{2\pi}\).
This law demonstrates that every irreducible physical tick requires a
non-punctual, complex relational cycle (\(2\pi\) internal phase closure
across spatial simultaneity span \(\mathcal{L}_{\mathrm{sim}}\))
balancing internal action-cadence conservation (\(\hbar c / 2\pi\))
against external boundary interaction (\(k_B T\)). We trace this unified
architecture across four physical scales (\emph{Star \(\to\) Atom
\(\to\) Electron \(\to\) Photon}), derive the exact \(1/4\)
Bekenstein--Hawking black hole entropy area law
(\(S = k_B A / 4\ell_P^2\)) from minimal clock-packing
(\(2\ell_P \times 2\ell_P\)) at the Quantum-Gravity Wall
(\(\Xi_{\mathrm{QG}} = 1\)), and apply the Horizon-Condensation Law to
five empirical domains: high-temperature superconductivity (\(T_c\)),
fundamental computational speed-cooling bounds (\(I_{\mathrm{sync}}\)),
the quark--gluon plasma shear-viscosity floor
(\(\eta/s \ge \hbar/4\pi k_B\)), the century-old Planck--Ott--Landsberg
relativistic temperature controversy, and sequential cosmological
freeze-out (\(A_2^1 \to A_3^1\)).

\begin{center}\rule{0.5\linewidth}{0.5pt}\end{center}

\subsection{1. Introduction: Reversing Kinematics and Thermodynamics
From Objective
Time}\label{introduction-reversing-kinematics-and-thermodynamics-from-objective-time}

Textbook physics treats space (\(\Delta x\)) and coordinate time
(\(\Delta t\)) as primary containers, defining speed
(\(v = \Delta x / \Delta t\)) as the trajectory of a zero-width point
moving across an empty stage and temperature (\(T\)) as the average
kinetic agitation of microscopic particles bouncing within that
container. Yet as established in the axiomatic foundations of the Theory
of Objective Time (\(\text{AEON-}T_{\mu\nu}\)) {[}1, 2{]}, \textbf{Being
is in itself and lasts} (\(A_s^v\)), whereas space and metric time
(\(g_{\mu\nu}\)) are never self-subsistent objects---they are secondary
operational survey maps earned by bounded, open-system physical
oscillators (\textbf{Centers of Resolution}, \(m > 0\)) phase-locking
within a finite \textbf{Operational Simultaneity Horizon}
(\(\mathcal{L}_{\mathrm{sim}}\)).

Once Objective Time---expressed both as raw un-clocked temporal
intensity (\(\rho_\tau > 0\)) and as discrete ordinal progression
(\(v \to v + 1\))---is placed at the foundation of physical reality, the
direction of explanation inverts. Neither speed nor temperature is a
secondary statistical accident; both emerge directly from how a
persistent entity partitions its finite temporal cadence between
\textbf{internal self-maintenance} and \textbf{external relational
interaction}.

\subsubsection{1.1 Speed as Coherent Cadence
Allocation}\label{speed-as-coherent-cadence-allocation}

An un-clocked excitation such as a massless photon (\(m = 0\)) lacks
Higgs chiral phase-locking (\(\psi_L \leftrightarrow \psi_R = 0\)) and
possesses zero internal rest-frame oscillation frequency
(\(\omega_C = 0\), \(d\tau = 0\)). Ruled purely by un-anchored temporal
cadence (\(\rho_\tau > 0\)), it allocates 100\% of its existence to
external spatial propagation at the local substrate cadence:

\[v_{\gamma} = \frac{c}{n_{\mathrm{eff}}(\mathbf{x})}\]

where
\(n_{\mathrm{eff}}(\mathbf{x}) = (dt_{\mathrm{pre}}(\mathbf{x})/dt_\infty)^{-1} \ge 1\)
is the local substrate metric impedance induced by surrounding
mass-anchored clocks {[}3{]}. Thus, \(c\) is not an arbitrary cosmic
speed limit imposed on matter; it is the fundamental dimensional
conversion factor linking raw temporal cadence to operational spatial
separation
(\(\Delta\ell = c\,\Delta\tau_{\mathrm{sig}} / n_{\mathrm{eff}}\)).

Conversely, a mass-anchored \textbf{Center of Resolution} (\(m > 0\))
locks raw cadence into an internal Compton phase-locking rhythm
\(\omega_C = m c^2 / \hbar\). At rest relative to an external survey
network (\(v_{\mathrm{cm}} = 0\), \(n_{\mathrm{eff}} = 1\)), 100\% of
its cadence sustains internal ordinal ticking
(\(dt_{\mathrm{pre}} = dt_\infty\)). Setting the clock into uniform
translational motion (\(0 < v_{\mathrm{cm}} < c\)) requires diverting a
fraction of that finite cadence budget into updating its external
spatial situating (\(s_1 \to s_2\)) across the substrate, reducing its
surviving internal tick rate by the Lorentz factor
\(\gamma = (1 - v_{\mathrm{cm}}^2/c^2)^{-1/2}\):

\[dt_{\mathrm{pre}} = \frac{dt}{n_{\mathrm{eff}}(\mathbf{x})}\sqrt{1 - \frac{v_{\mathrm{cm}}^2}{c^2}}\]

Kinematic time dilation is therefore the direct operational cost of
coherent spatial re-situating. Reaching \(v_{\mathrm{cm}} \to c\) is
ontologically forbidden for any bound clock (\(m > 0\)) because
diverting 100\% of cadence to external displacement leaves zero cycles
for internal chiral phase-locking, dissolving the Center of Resolution
itself.

\subsubsection{1.2 Temperature as Incoherent Multi-Node Condensation
Rate}\label{temperature-as-incoherent-multi-node-condensation-rate}

Whereas speed measures \textbf{coherent, directional} cadence
expenditure along a unified trajectory, \textbf{temperature (\(T\))}
measures the \textbf{incoherent, stochastic rate} at which open-system
boundary interactions (\(\Delta\Omega > 0\)) precipitate irreversible
\textbf{Two-Stage Condensations} (\(\Phi_\tau(\Omega)\)) across an
ensemble of coupled Centers of Resolution {[}4{]}.

Within a thermal bath at temperature \(T\), environmental interactions
inject random phase perturbations at the characteristic thermal
condensation frequency:

\[\omega_{\mathrm{th}} = \frac{k_B T}{\hbar}\]

\begin{itemize}
\tightlist
\item
  \textbf{The Low-Temperature Floor (\(T \to 0\)):} Stochastic boundary
  condensation noise vanishes (\(\Gamma_{\mathrm{tot}} \to 0\)),
  permitting the \textbf{Operational Simultaneity Horizon}
  (\(\mathcal{L}_{\mathrm{sim}} = c\,n_{\mathrm{eff}}^{-1}/\Gamma_{\mathrm{tot}}\))
  to expand across macroscopic distances so multiple oscillators
  phase-lock into a single unified quantum domain.
\item
  \textbf{The High-Temperature Meltdown Ceiling
  (\(\Xi_{\mathrm{QG}} \to 1\)):} As \(T\) rises, stochastic
  condensation ticks interrupt internal phase-locking ever more rapidly.
  When thermal agitation energy approaches the bound oscillator's
  rest-mass locking scale (\(k_B T \sim m c^2\)), external boundary
  interactions shatter internal chiral coherence within a single tick,
  melting bound clocks back into un-clocked conformal radiation
  (\(A_3^1 \to A_2^1\)) {[}5{]}.
\end{itemize}

\begin{center}\rule{0.5\linewidth}{0.5pt}\end{center}

\subsection{2. The Cadence Partition Identity and the Weyl Scale
Elevator}\label{the-cadence-partition-identity-and-the-weyl-scale-elevator}

Combining gravitational pacing impedance (\(n_{\mathrm{eff}}\)),
coherent translational speed (\(v_{\mathrm{cm}}/c\)), and incoherent
thermal condensation load (\(k_B T_{\mathrm{eff}} / m c^2\)) yields the
\textbf{Cadence Partition Identity} governing the surviving internal
proper-time rate of any bound Center of Resolution (\(m > 0\)):

\[\left(\frac{dt_{\mathrm{pre}}}{dt_\infty}\right)^2 = \frac{1}{n_{\mathrm{eff}}^2(\mathbf{x})} \left(1 - \frac{v_{\mathrm{cm}}^2}{c^2}\right) \left(1 - \frac{k_B T_{\mathrm{eff}}}{m c^2}\right)\]

This identity reveals that \textbf{ordered speed} and \textbf{disordered
temperature} enter the survival budget of a physical clock
symmetrically. Pushing either \(v_{\mathrm{cm}} \to c\) or
\(k_B T_{\mathrm{eff}} \to m c^2\) drives the surviving internal clock
cadence \(dt_{\mathrm{pre}} \to 0\), terminating at the exact same
operational boundary---the \textbf{Quantum-Gravity Wall}
(\(\Xi_{\mathrm{QG}} = 1\))---where bound clocks cease to persist
{[}6{]}.

\subsubsection{2.1 Invariance Across the Weyl Scale
Elevator}\label{invariance-across-the-weyl-scale-elevator}

Under a global \textbf{Weyl scale transformation} moving up or down
energy/frequency regimes (\(E \to \lambda E\),
\(\omega \to \lambda \omega\)), characteristic mass-energy scales,
thermal frequencies, and horizon lengths transform covariantly:

\[m c^2 \longrightarrow \lambda\,m c^2, \qquad k_B T \longrightarrow \lambda\,k_B T, \qquad \mathcal{L}_{\mathrm{sim}} \longrightarrow \lambda^{-1}\,\mathcal{L}_{\mathrm{sim}}\]

Consequently, both the velocity ratio (\(v_{\mathrm{cm}}/c\)) and the
thermal-to-rest-cadence ratio (\(k_B T_{\mathrm{eff}} / m c^2\)) remain
\textbf{strictly Weyl-invariant} across every rung of the scale
hierarchy:

{\def\LTcaptype{none} % do not increment counter
\begin{longtable}[]{@{}
  >{\raggedright\arraybackslash}p{(\linewidth - 8\tabcolsep) * \real{0.2000}}
  >{\raggedright\arraybackslash}p{(\linewidth - 8\tabcolsep) * \real{0.2000}}
  >{\raggedright\arraybackslash}p{(\linewidth - 8\tabcolsep) * \real{0.2000}}
  >{\raggedright\arraybackslash}p{(\linewidth - 8\tabcolsep) * \real{0.2000}}
  >{\raggedright\arraybackslash}p{(\linewidth - 8\tabcolsep) * \real{0.2000}}@{}}
\toprule\noalign{}
\begin{minipage}[b]{\linewidth}\raggedright
Weyl Scale Regime
\end{minipage} & \begin{minipage}[b]{\linewidth}\raggedright
Characteristic Cadence Floor
\end{minipage} & \begin{minipage}[b]{\linewidth}\raggedright
Coherent Speed Mode (\(v/c\))
\end{minipage} & \begin{minipage}[b]{\linewidth}\raggedright
Incoherent Thermal Mode (\(T\))
\end{minipage} & \begin{minipage}[b]{\linewidth}\raggedright
Unified Operational Behavior
\end{minipage} \\
\midrule\noalign{}
\endhead
\bottomrule\noalign{}
\endlastfoot
\textbf{Pre-Metric Conformal Pool (\(A_2^1\))} & Un-clocked
\(\rho_\tau > 0\) (\(\omega_C = 0\)) & Pure substrate cadence
(\(v = c\)) & Scale-free conformal bath (\(T \cdot a = \text{const}\)) &
Speed and temperature merge into a single massless radiative flux. \\
\textbf{Electroweak Nucleation (\(\sim 100\ \text{GeV}\))} & Chiral
lock-in (\(\omega_H \sim v_H/\hbar\)) & Sub-luminal inertia (\(v < c\))
& Symmetry-breaking threshold (\(T_{\mathrm{EW}}\)) & Unified radiation
bifurcates into massive oscillators and thermal bath. \\
\textbf{Quantum \& Atomic Regime (\(\text{eV} - \text{MeV}\))} & Compton
rhythm (\(\omega_C = mc^2/\hbar\)) & Wave-packet drift (\(v \ll c\)) &
Coherence gate (\(\mathcal{L}_{\mathrm{sim}} \propto 1/T\)) &
Acceleration converts ordered speed into Unruh horizon heat
(\(T \propto a\)). \\
\textbf{Classical Macroscopic Regime (\(300\ \text{K}\))} & Low-shear
valley (\(\Xi_{\mathrm{QG}} \ll 1\)) & Collective drift
(\(v_{\mathrm{cm}}\)) & Kinetic micro-cadence
(\(\frac{1}{2}m\langle v_{\mathrm{th}}^2\rangle = \frac{3}{2}k_B T\)) &
Macroscopic temperature is revealed as incoherent microscopic speed. \\
\textbf{Quantum-Gravity Wall (\(\Xi_{\mathrm{QG}} \to 1\))} & Breakdown
ceiling (\(\omega_{\mathrm{tidal}} \to \omega_C\)) & Horizon escape
limit (\(v_{\mathrm{esc}} \to c/n_{\mathrm{eff}}\)) & Hagedorn / Hawking
ceiling (\(k_B T \to mc^2/2\pi\)) & Bound clocks dissolve back into pure
un-clocked temporal intensity (\(\rho_\tau\)). \\
\end{longtable}
}

\begin{center}\rule{0.5\linewidth}{0.5pt}\end{center}

\subsection{3. Acceleration as Horizon Truncation and the
Horizon-Condensation
Law}\label{acceleration-as-horizon-truncation-and-the-horizon-condensation-law}

What physically links coherent speed changes to incoherent thermal
condensation? The bridge is \textbf{acceleration} (\(a = dv/dt\))---and,
by the Equivalence Principle (\(m_i \equiv m_g\)) derived from chiral
phase-locking {[}3{]}, \textbf{gravitational pacing gradients}
(\(\nabla \rho_\tau\)).

\subsubsection{\texorpdfstring{3.1 Horizon Truncation
(\(a \to T\))}{3.1 Horizon Truncation (a \textbackslash to T)}}\label{horizon-truncation-a-to-t}

When a bound Center of Resolution undergoes uniform proper acceleration
\(a\), round-trip phase-locking signals (\(\Delta\tau_{\mathrm{sig}}\))
cannot return from distances exceeding the causal Rindler horizon
\(L_a = c^2 / a\). Acceleration therefore physically clamps the clock's
\textbf{Operational Simultaneity Horizon} to:

\[\mathcal{L}_{\mathrm{sim}}(a) = \frac{c^2}{a}\]

In an inertial frame (\(a = 0\),
\(\mathcal{L}_{\mathrm{sim}} \to \infty\)), vacuum modes within
\textbf{Functional Infinity} (\(\Omega\)) remain phase-correlated across
the horizon and cancel cleanly without triggering condensation. Clamping
\(\mathcal{L}_{\mathrm{sim}}(a) = c^2/a\) severs vacuum phase
correlations across the boundary, injecting uncorrelated open-system
entropy noise (\(\Delta\Omega > 0\)) into the accelerating clock. This
forces \textbf{Two-Stage Condensation} (\(\Phi_\tau\)) to fire at the
stochastic rate \(\Gamma_{\mathrm{cond}} = a / (2\pi c)\), precipitating
real thermal energy at the Unruh temperature:

\[k_B T_{\mathrm{Unruh}} = \hbar\,\Gamma_{\mathrm{cond}} = \frac{\hbar a}{2\pi c}\]

\subsubsection{\texorpdfstring{3.2 Derivation of the
Horizon-Condensation Constant
(\(\mathcal{C}_\tau = 1/2\pi\))}{3.2 Derivation of the Horizon-Condensation Constant (\textbackslash mathcal\{C\}\_\textbackslash tau = 1/2\textbackslash pi)}}\label{derivation-of-the-horizon-condensation-constant-mathcalc_tau-12pi}

Eliminating the acceleration \(a\) between the truncated simultaneity
horizon \(\mathcal{L}_{\mathrm{sim}} = c^2 / a\) and the precipitated
thermal energy \(k_B T = \hbar a / (2\pi c)\) causes both \(a\) and the
Weyl scale factor \(\lambda\) to cancel identically, yielding a
universal, dimensionless constant:

\[\mathcal{C}_\tau \equiv \frac{k_B T \cdot \mathcal{L}_{\mathrm{sim}}}{\hbar c} = \frac{1}{2\pi} \approx 0.159155\]

We designate this relation \textbf{The Horizon-Condensation Law}, and
\(\mathcal{C}_\tau = 1/2\pi\) as the \textbf{Chronometric
Horizon-Condensation Constant}.

\subsubsection{3.3 The Complex Structure of a Situated Moment:
Conservation and
Interaction}\label{the-complex-structure-of-a-situated-moment-conservation-and-interaction}

Written in energy-action balance form, the Horizon-Condensation Law
exposes the internal anatomy of an irreducible physical tick
(\(A_s^v \to A_{s'}^{v+1}\)):

\[\underbrace{k_B T}_{\substack{\text{Boundary Interaction} \\ (\text{Open-System Coupling } \Delta\Omega > 0)}} = \underbrace{\frac{1}{2\pi}\,\frac{\hbar c}{\mathcal{L}_{\mathrm{sim}}}}_{\substack{\text{Internal Conservation} \\ (\text{One } 2\pi \text{ Action-Cadence Cycle})}}\]

In textbook point-continuum physics, a ``moment'' is idealized as a
zero-width slice (\(t = t_0\)) possessing neither thickness nor spatial
breadth. The Horizon-Condensation Law proves that a physical moment is
necessarily \textbf{complex}, spanning two structural dimensions and
governed by two non-negotiable physical laws:

\begin{enumerate}
\def\labelenumi{\arabic{enumi}.}
\tightlist
\item
  \textbf{Internal Cyclic Depth (\(2\pi\) Radians) \& Conservation:} A
  zero-width point cannot register change. To advance its ordinal index
  (\(v \to v + 1\)), a Center of Resolution must execute \textbf{one
  complete \(2\pi\) phase-locking cycle} across its simultaneity span
  \(\mathcal{L}_{\mathrm{sim}}\). The factor
  \(\hbar c / (2\pi \mathcal{L}_{\mathrm{sim}})\) enforces strict
  \textbf{conservation} of action (\(\hbar\)), causal cadence (\(c\)),
  and conserved charges across the transition, forbidding arbitrary
  discontinuities.
\item
  \textbf{Spatial Simultaneity Breadth (\(\mathcal{L}_{\mathrm{sim}}\))
  \& Interaction:} A completely isolated system (\(\Delta\Omega = 0\),
  \(T = 0\) with \(\mathcal{L}_{\mathrm{sim}} \to \infty\)) remains
  suspended inside uncondensed wave potential (\(\Omega\)) without
  precipitating a classical record. Real \textbf{boundary interaction}
  (\(k_B T > 0\) at finite \(\mathcal{L}_{\mathrm{sim}}\)) supplies the
  exact condensation threshold required to trigger \(\Phi_\tau\) and
  lock in the successor state \(A_{s'}^{v+1}\).
\end{enumerate}

Setting \(\mathcal{L}_{\mathrm{sim}}\) equal to a bound particle's own
internal Compton coherence width
(\(\mathcal{L}_{\mathrm{sim}} = \lambda_C = \hbar / mc\), where
\(\Xi_{\mathrm{QG}} = 1\)) immediately yields the universal
\textbf{Compton--Unruh Thermal Meltdown Ceiling} for any bound clock of
mass \(m\):

\[\frac{k_B T_{\mathrm{crit}}}{m c^2} = \mathcal{C}_\tau = \frac{1}{2\pi} \approx 0.1592\]

A bound Center of Resolution cannot survive up to \(k_B T = m c^2\); as
soon as stochastic horizon condensation reaches
\(1/2\pi \approx 15.92\%\) of its rest-mass energy, internal chiral
phase-locking snaps and pair-production dissolves the isolated
oscillator.

\begin{center}\rule{0.5\linewidth}{0.5pt}\end{center}

\subsection{\texorpdfstring{4. Multi-Scale Anatomy of a Single Situated
Moment: Star \(\to\) Atom \(\to\) Electron \(\to\)
Photon}{4. Multi-Scale Anatomy of a Single Situated Moment: Star \textbackslash to Atom \textbackslash to Electron \textbackslash to Photon}}\label{multi-scale-anatomy-of-a-single-situated-moment-star-to-atom-to-electron-to-photon}

Because \(\mathcal{C}_\tau = 1/2\pi\) is Weyl-invariant, the exact same
operational law governs reality simultaneously across all zoom levels at
a single situated moment, sliding smoothly along the Weyl Scale Elevator
as \(\mathcal{L}_{\mathrm{sim}}\) narrows:

\begin{enumerate}
\def\labelenumi{\arabic{enumi}.}
\tightlist
\item
  \textbf{Macroscopic Star (e.g., The Sun,
  \(R \approx 6.96 \times 10^8\ \text{m}\)):}\\
  A star is not a single coherent quantum clock; it is a composite
  classical congregation of \(\sim 10^{57}\) decoupled Centers of
  Resolution whose local simultaneity horizons
  (\(\mathcal{L}_{\mathrm{sim}} \sim 10^{-11}\ \text{m}\) in the hot
  plasma core) are orders of magnitude smaller than the stellar radius
  (\(L_{\mathrm{star}} \gg \mathcal{L}_{\mathrm{sim}}\)). Collectively,
  their concentrated rest-mass cadence carves a deep chronometric well
  (\(\Delta n_{\mathrm{eff}} \approx 2.12 \times 10^{-6}\) at the
  photosphere). Globally, stellar hydrostatic equilibrium represents the
  exact balance between outward stochastic thermal condensation pressure
  (\(T_{\mathrm{core}} \approx 1.5 \times 10^7\ \text{K}\)) and inward
  chronometric pacing gradients (\(\nabla \rho_\tau\)).
\item
  \textbf{Bound Atom (e.g., Hydrogen or Cesium,
  \(R \sim 10^{-10}\ \text{m}\)):}\\
  Zooming into a tenuous stellar corona or an optical-lattice vacuum
  chamber, electromagnetic phase-locking binds nucleus and electron
  cloud within a single shared \textbf{Operational Simultaneity Horizon}
  (\(\mathcal{L}_{\mathrm{sim}} \ge R_{\mathrm{atom}}\)). Here, one
  atomic tick (\(v \to v + 1\)) corresponds to a full \(2\pi\)
  standing-wave orbital phase closure across
  \(\sim 10^{-10}\ \text{m}\). So long as ambient thermal condensation
  stays below \(k_B T \ll \hbar c / (2\pi \mathcal{L}_{\mathrm{sim}})\),
  the atom persists as one unified composite clock; exceeding that bound
  triggers ionization via stochastic condensation (\(\Phi_\tau\)).
\item
  \textbf{Elementary Electron
  (\(\lambda_C \approx 2.43 \times 10^{-12}\ \text{m}\)):}\\
  Zooming past the atomic shell to a solitary bound electron, Higgs
  chiral phase-locking (\(\psi_L \leftrightarrow \psi_R\)) anchors raw
  cadence into an irreducible internal Compton rhythm
  (\(\omega_C = m_e c^2 / \hbar \approx 7.76 \times 10^{20}\ \text{rad/s}\)).
  Between boundary couplings, its spatial situating opens into
  \textbf{Functional Infinity} (\(\Omega\)) across its Compton coherence
  envelope; upon boundary interaction (\(\Delta\Omega > 0\)), one
  \(2\pi\) cycle precipitates a localized particle event
  (\(A_s^v \to A_{s'}^{v+1}\)). Its solitary phase-lock remains intact
  until horizon temperature reaches
  \(k_B T_{\mathrm{crit}} = m_e c^2 / (2\pi) \approx 81.3\ \text{keV}\)
  (\(T_{\mathrm{crit}} \approx 9.44 \times 10^8\ \text{K}\)), where
  spontaneous \(e^- e^+\) pair condensation overwhelms individual
  identity.
\item
  \textbf{Massless Photon (\(m = 0\), \(\omega_C = 0\)):}\\
  Zooming all the way to an un-clocked photon in transit, chiral
  phase-locking vanishes (\(m = 0\)) and no internal rest-frame clock
  exists (\(d\tau = 0\)). Carrying pure \textbf{Objective Temporal
  Intensity} (\(\rho_\tau > 0\)) and conserved energy-momentum
  (\(\hbar\omega\)), the photon allocates 100\% of its cadence to
  external propagation (\(v = c/n_{\mathrm{eff}}\)) along the star's
  refractive gradient, remaining uncondensed until it couples to a bound
  Center of Resolution (an atomic electron in a retina or photodetector)
  and enters a complex \(2\pi\) moment.
\end{enumerate}

\begin{center}\rule{0.5\linewidth}{0.5pt}\end{center}

\subsection{5. Derivation of the Bekenstein--Hawking Horizon Area Law
From Clock
Packing}\label{derivation-of-the-bekensteinhawking-horizon-area-law-from-clock-packing}

Applied to the gravitational horizon of a Schwarzschild black hole of
mass \(M\), Schwarzschild radius \(R_s = 2GM/c^2\), surface area
\(A = 4\pi R_s^2\), and surface gravity
\(g_s = c^4/(4GM) = c^2/(2R_s)\), the Horizon-Condensation Law derives
the exact \(1/4\) Bekenstein--Hawking entropy formula without Euclidean
imaginary-time analytic continuation or higher-dimensional D-brane
counting {[}7, 8{]}:

\begin{enumerate}
\def\labelenumi{\arabic{enumi}.}
\item
  \textbf{Horizon Simultaneity Span:} A stationary boundary oscillator
  held at the Quantum-Gravity Wall (\(\Xi_{\mathrm{QG}} = 1\))
  experiences surface acceleration \(a = g_s = c^2/(2R_s)\), clamping
  its local Operational Simultaneity Horizon to the horizon diameter:

  \[\mathcal{L}_{\mathrm{sim}} = \frac{c^2}{g_s} = 2 R_s\]
\item
  \textbf{Hawking Condensation Temperature:} Substituting
  \(\mathcal{L}_{\mathrm{sim}} = 2 R_s\) directly into
  \(\mathcal{C}_\tau = \frac{k_B T \cdot \mathcal{L}_{\mathrm{sim}}}{\hbar c} = \frac{1}{2\pi}\)
  yields:

  \[k_B T_H = \left(\frac{1}{2\pi}\right)\frac{\hbar c}{2 R_s} = \frac{\hbar c}{4\pi R_s} = \frac{\hbar c^3}{8\pi G M}\]

  Here \(\mathcal{C}_\tau = 1/2\pi\) supplies the exact \(2\pi\)
  phase-cycle factor while round-trip horizon diameter
  \(\mathcal{L}_{\mathrm{sim}} = 2R_s\) supplies the factor of \(2\).
\item
  \textbf{Area Quantization From Minimal Clock Packing:} Applying
  \(dE = T_H\,dS\) for an incremental mass-energy condensation
  \(dE = c^2\,dM\), differentiating \(A = 16\pi G^2 M^2 / c^4\) to
  obtain \(8\pi G M\,dM = (c^4/4G)\,dA\), and substituting the Planck
  area \(\ell_P^2 = \hbar G / c^3\) gives:

  \[dS = \frac{c^2\,dM}{T_H} = k_B\,\frac{8\pi G M}{\hbar c}\,dM = k_B\,\frac{c^3}{4\hbar G}\,dA = k_B\,\frac{dA}{4\ell_P^2} \quad \Longrightarrow \quad S_{\mathrm{BH}} = k_B\,\frac{A}{4\ell_P^2}\]
\end{enumerate}

The famous denominator \textbf{\(4\)} acquires a direct operational
meaning: because internal singularity collapse (\(r = 0\)) is excised by
\(\mathcal{P}_{\mathrm{QG}}(\Xi_{\mathrm{QG}}) \to 0\) at
\(\Xi_{\mathrm{QG}} = 1\) {[}3{]}, all physical microstates reside on
the boundary impedance shell, tiled by elementary Centers of Resolution
each occupying a minimum two-way phase-locking cell of diameter
\(2\ell_P\) (\(2\ell_P \times 2\ell_P = 4\ell_P^2\)) and condensing at
one bit (\(k_B\)) per \(2\pi\) cycle.

\begin{center}\rule{0.5\linewidth}{0.5pt}\end{center}

\subsection{6. Applications of the Horizon-Condensation Law Across Five
Physical
Regimes}\label{applications-of-the-horizon-condensation-law-across-five-physical-regimes}

Because
\(\mathcal{C}_\tau = \frac{k_B T \cdot \mathcal{L}_{\mathrm{sim}}}{\hbar c} = \frac{1}{2\pi}\)
governs every interface between coherent phase-locking
(\(\mathcal{L}_{\mathrm{sim}}\)) and stochastic thermal condensation
(\(k_B T\)), it yields quantitative predictions and structural
resolutions across five distinct domains of physics:

\subsubsection{\texorpdfstring{6.1 Macroscopic Quantum Coherence and the
Superconducting Critical Ceiling
(\(T_c\))}{6.1 Macroscopic Quantum Coherence and the Superconducting Critical Ceiling (T\_c)}}\label{macroscopic-quantum-coherence-and-the-superconducting-critical-ceiling-t_c}

In a normal metal, conduction electrons act as independent Centers of
Resolution whose individual simultaneity horizons are shorter than the
inter-electron spacing, so lattice interactions continuously trigger
irreversible condensations (\(\Phi_\tau\)) manifested as electrical
resistance. Below the critical temperature \(T_c\), pairing interactions
lock Cooper pairs across coherence length \(\xi_0\) into \textbf{one
unified macroscopic Operational Simultaneity Horizon}
(\(\mathcal{L}_{\mathrm{sim}} \ge \xi_0\)). Applying the
Horizon-Condensation Law with characteristic carrier phase-velocity
\(v_F\) bounds the maximum thermal condensation load the coherent state
can sustain before localized ticks shatter macroscopic phase-lock:

\[k_B T_c \le \mathcal{C}_\tau\,\frac{\hbar v_F}{\xi_0} = \frac{\hbar v_F}{2\pi\,\xi_0}\]

This reproduces the structural scaling of the
Bardeen--Cooper--Schrieffer (BCS) coherence relation
(\(\xi_0 \sim \hbar v_F / \pi\Delta_0\)) {[}9{]}, showing that raising
\(T_c\) requires engineering denser, shorter-range pairing horizons
(\(\xi_0 \to \text{nm}\)) capable of outpacing stochastic condensation
ticks.

\subsubsection{\texorpdfstring{6.2 Quantum Computation Scaling and the
Synchronization Tax
(\(I_{\mathrm{sync}}\))}{6.2 Quantum Computation Scaling and the Synchronization Tax (I\_\{\textbackslash mathrm\{sync\}\})}}\label{quantum-computation-scaling-and-the-synchronization-tax-i_mathrmsync}

Keeping an \(N\)-qubit quantum processor of physical span
\(L_{\mathrm{chip}}\) suspended inside uncondensed \textbf{Functional
Infinity} (\(\Omega\)) without premature decoherence requires
maintaining a single shared simultaneity horizon
\(\mathcal{L}_{\mathrm{sim}} \ge L_{\mathrm{chip}}\) across all coupled
nodes. By the Horizon-Condensation Law, operating temperature is bounded
strictly by:

\[T_{\mathrm{op}} \ll \frac{\hbar c}{2\pi k_B L_{\mathrm{chip}}\,\mathcal{N}_{\mathrm{couple}}(N)}\]

where \(\mathcal{N}_{\mathrm{couple}}(N)\) accounts for the cumulative
boundary surface exposed to environmental tick leakage
(\(\Delta\Omega > 0\)). Expanding chip footprint (\(L_{\mathrm{chip}}\))
or qubit connectivity drives the \textbf{Synchronization Tax}
(\(I_{\mathrm{sync}}\)) upward, unifying Landauer's thermodynamic
erasure limit (\(k_B T \ln 2\)) {[}10{]} and the Margolus--Levitin
quantum speed limit {[}11{]} into a single operational constraint per
\(2\pi\) logical cycle.

\subsubsection{\texorpdfstring{6.3 Quark--Gluon Plasma and the Universal
Shear-Viscosity Bound
(\(\eta/s \ge \hbar / 4\pi k_B\))}{6.3 Quark--Gluon Plasma and the Universal Shear-Viscosity Bound (\textbackslash eta/s \textbackslash ge \textbackslash hbar / 4\textbackslash pi k\_B)}}\label{quarkgluon-plasma-and-the-universal-shear-viscosity-bound-etas-ge-hbar-4pi-k_b}

In ultra-relativistic heavy-ion collisions at RHIC and the LHC, hadronic
matter melts above \(T_{\mathrm{QCD}} \approx 170\ \text{MeV}\) into a
strongly coupled quark--gluon plasma (QGP) exhibiting a near-perfect
fluid shear-viscosity-to-entropy-density ratio bounded from below by the
Kovtun--Son--Starinets (KSS) limit {[}12{]}:

\[\frac{\eta}{s} \ge \frac{\hbar}{4\pi k_B}\]

Whereas standard derivations invoke five-dimensional Anti-de
Sitter/Conformal Field Theory (AdS/CFT) string dualities, the
Horizon-Condensation Law derives this exact bound directly in four
operational dimensions. Notice that
\(\frac{\hbar}{4\pi k_B} = \frac{1}{2}\,\mathcal{C}_\tau\,\frac{\hbar}{k_B}\).
At horizon saturation---where thermal condensation energy (\(k_B T\))
shrinks \(\mathcal{L}_{\mathrm{sim}}\) down to the quantum mean free
path of the excitations---every single \(2\pi\) cycle precipitates a
boundary condensation carrying one quantum of entropy
(\(s \sim n\,2\pi k_B\)) and one half-cycle (\(\pi\)) of transverse
momentum shear (\(\eta \sim n\,\hbar / 2\)), immediately fixing
\(\eta / s = \hbar / (4\pi k_B)\).

\subsubsection{6.4 Resolution of the Planck--Ott--Landsberg Relativistic
Temperature
Controversy}\label{resolution-of-the-planckottlandsberg-relativistic-temperature-controversy}

For over a century, relativistic thermodynamics has been divided over
how the temperature \(T\) of a moving body transforms under a Lorentz
boost of speed \(v\) {[}13--15{]}: Einstein and Planck derived
\(T' = T/\gamma\) (moving bodies appear colder), Ott derived
\(T' = \gamma T\) (moving bodies appear hotter), and Landsberg argued
for \(T' = T\) (Lorentz invariance).

The Cadence Partition Identity and the Horizon-Condensation Law show
that each formulation isolates a different operational variable while
leaving \(\mathcal{C}_\tau\) strictly invariant: 1. \textbf{Internal
Comoving Tick Rate (Einstein--Planck):} Because translational speed
consumes part of the ensemble's cadence budget
(\(dt_{\mathrm{pre}} = dt/\gamma\)), an internal comoving thermometer
registers a slower proper condensation rate per coordinate second
(\(T_{\mathrm{rate}}' = T/\gamma\)). 2. \textbf{Longitudinal Horizon
Contraction (Ott):} Simultaneously, relativistic motion contracts the
longitudinal \textbf{Operational Simultaneity Horizon} along the boost
direction
(\(\mathcal{L}_{\mathrm{sim},\parallel}' = \mathcal{L}_{\mathrm{sim}}/\gamma\)).
A stationary external detector coupling across that compressed horizon
measures an amplified boundary condensation flux
\(k_B T_{\mathrm{det}}' = \hbar c / [2\pi(\mathcal{L}_{\mathrm{sim}}/\gamma)] = \gamma\,k_B T\).
3. \textbf{True Lorentz \& Weyl Invariant:} Neither \(T\) alone nor
\(\mathcal{L}_{\mathrm{sim}}\) alone is frame-invariant; what remains
invariant across every inertial frame is their dimensionless product:

\[\mathcal{C}_\tau' = \frac{k_B T_{\mathrm{det}}' \cdot \mathcal{L}_{\mathrm{sim},\parallel}'}{\hbar c} = \frac{(\gamma\,k_B T)(\mathcal{L}_{\mathrm{sim}}/\gamma)}{\hbar c} = \frac{1}{2\pi} = \mathcal{C}_\tau\]

\subsubsection{\texorpdfstring{6.5 Cosmological Phase Transitions and
Sequential Mass Freeze-Out
(\(A_2^1 \to A_3^1\))}{6.5 Cosmological Phase Transitions and Sequential Mass Freeze-Out (A\_2\^{}1 \textbackslash to A\_3\^{}1)}}\label{cosmological-phase-transitions-and-sequential-mass-freeze-out-a_21-to-a_31}

As the early universe expanded after the pre-metric crossover
(\(A_2^1 \to A_3^1\)), ambient radiation cooled inversely with the
cosmic scale factor (\(T(a) \propto a^{-1}\)), driving a continuous
widening of the cosmic \textbf{Operational Simultaneity Horizon}:

\[\mathcal{L}_{\mathrm{sim}}(a) = \frac{\hbar c}{2\pi\,k_B T(a)} \propto a\]

A bound \textbf{Center of Resolution} of mass \(m\) and internal Compton
width \(\lambda_C(m) = \hbar / mc\) can stabilize against thermal
melting (\(\Phi_\tau\) disruption) only once the expanding simultaneity
horizon encloses its internal \(2\pi\) phase-locking cycle
(\(\mathcal{L}_{\mathrm{sim}}(a) \ge \lambda_C(m)\)). This fixes the
sequential cosmological freeze-out temperature for every particle
species in cosmic history:

\[k_B T_{\mathrm{freeze}}(m) = \mathcal{C}_\tau\,m c^2 = \frac{m c^2}{2\pi}\]

Cosmic structure thus emerges as a progressive descent down the Weyl
Scale Elevator: as \(\mathcal{L}_{\mathrm{sim}}(a)\) expands, reality
sequentially crosses the stabilization thresholds for electroweak bosons
(\(\sim 100\ \text{GeV}\)), QCD hadrons (\(\sim 170\ \text{MeV}\)),
electron--positron annihilation
(\(k_B T_{\mathrm{crit}} \approx 81.3\ \text{keV}\)), and neutral
hydrogen recombination (\(\sim 0.3\ \text{eV}\)).

\begin{center}\rule{0.5\linewidth}{0.5pt}\end{center}

\subsection{7. Conclusion}\label{conclusion}

By grounding physical reality in Objective Time (\(\rho_\tau > 0\)) and
bounded, open-system physical clocks (\(A_s^v\)), \textbf{Speed} and
\textbf{Temperature} are unified as the coherent and incoherent modes of
cadence expenditure across the Weyl Scale Elevator. From the causal
truncation of simultaneity horizons under acceleration
(\(\mathcal{L}_{\mathrm{sim}} = c^2/a\)) emerges \textbf{The
Horizon-Condensation Law},
\(\mathcal{C}_\tau = \frac{k_B T \cdot \mathcal{L}_{\mathrm{sim}}}{\hbar c} = \frac{1}{2\pi}\).
This Weyl- and Lorentz-invariant constant demonstrates that every
irreducible moment of physical reality is a complex
\(2\pi \times \mathcal{L}_{\mathrm{sim}}\) relational cycle balancing
internal action-cadence conservation against external boundary
interaction---simultaneously unlocking the \(1/4\) Bekenstein--Hawking
black hole entropy law, superconducting coherence ceilings, quantum
processor cooling limits, the quark--gluon plasma viscosity bound,
relativistic temperature covariance, and sequential cosmological
freeze-out.

\begin{center}\rule{0.5\linewidth}{0.5pt}\end{center}

\subsection{References}\label{references}

\begin{enumerate}
\def\labelenumi{\arabic{enumi}.}
\tightlist
\item
  Pimenta Damásio, L. (2026). \emph{The Theory of Objective Time:
  Axiomatic Foundations of the \(\text{AEON-}T_{\mu\nu}\) Framework:
  Ontological Postulates, Operational Metrics, and Fundamental Laws of
  Temporal Actualization}. Zenodo.
  https://doi.org/10.5281/zenodo.24392857
\item
  Pimenta Damásio, L. (2026). \emph{A Post-Relativistic Operational
  Process Ontology: Situating the Theory of Objective Time
  (\(\text{AEON-}T_{\mu\nu}\)) Within the Philosophical and Physical
  Tradition}. Zenodo. https://doi.org/10.5281/zenodo.24442612
\item
  Pimenta Damásio, L. (2026). \emph{From Chiral Phase-Locking to
  Chronometric Gravity: Higgs Mass Nucleation, Centers of Resolution,
  and Falsifiable Quantum-Gravity Signatures in the Theory of Objective
  Time}. Zenodo. https://doi.org/10.5281/zenodo.24424098
\item
  Pimenta Damásio, L. (2026). \emph{From Pre-Quantitative Infinities to
  Diachronic Logic: Pre-Axiomatic Foundations of the Theory of Objective
  Time}. Zenodo. https://doi.org/10.5281/zenodo.24365866
\item
  Pimenta Damásio, L. (2026). \emph{From Pre-Metric Cadence to Earned
  Cosmic Time: Tracing the Big Bang Across the Aeon Crossover
  (\(A = A_1^1 \to A_2^1 \to A_3^1\))}. Zenodo.
  https://doi.org/10.5281/zenodo.24420549
\item
  Pimenta Damásio, L. (2026). \emph{Where Equations Break Down: Four
  Pathologies of Modern Physics Resolved Through the Theory of Objective
  Time}. Zenodo. https://doi.org/10.5281/zenodo.24421529
\item
  Bekenstein, J. D. (1973). Black holes and entropy. \emph{Physical
  Review D}, 7(8), 2333--2346.
\item
  Hawking, S. W. (1975). Particle creation by black holes.
  \emph{Communications in Mathematical Physics}, 43(3), 199--220.
\item
  Bardeen, J., Cooper, L. N., \& Schrieffer, J. R. (1957). Theory of
  superconductivity. \emph{Physical Review}, 108(5), 1175--1204.
\item
  Landauer, R. (1961). Irreversibility and heat generation in the
  computing process. \emph{IBM Journal of Research and Development},
  5(3), 183--191.
\item
  Margolus, N., \& Levitin, L. B. (1998). The maximum speed of dynamical
  evolution. \emph{Physica D: Nonlinear Phenomena}, 120(1--2), 188--195.
\item
  Kovtun, P., Son, D. T., \& Starinets, A. O. (2005). Viscosity in
  strongly interacting quantum field theories from black hole physics.
  \emph{Physical Review Letters}, 94(11), 111601.
\item
  Planck, M. (1908). Zur Dynamik bewegter Systeme. \emph{Annalen der
  Physik}, 331(6), 1--34.
\item
  Ott, H. (1963). Lorentz-Transformation der Wärme und der Temperatur.
  \emph{Zeitschrift für Physik}, 175(1), 70--104.
\item
  Landsberg, P. T. (1966). Does a moving body appear cool?
  \emph{Nature}, 212(5062), 571--572.
\end{enumerate}

\begin{center}\rule{0.5\linewidth}{0.5pt}\end{center}

\subsection{AI Assistance Disclosure}\label{ai-assistance-disclosure}

In accordance with open-science transparency standards, the author
discloses that an AI research assistant (Google Gemini) was utilized
during the preparation of this manuscript to verify dimensional
consistency across Weyl scale transformations, structure comparative
tables, and format mathematical derivations in LaTeX/Markdown. All
ontological postulates, conceptual definitions, physical
interpretations, and final editorial decisions are solely those of the
author.

\end{document}
